% 07 — Modus Ponens : รูปแบบกฎการอนุมาน (premises เหนือเส้น, ข้อสรุปใต้เส้น) คู่กับตัวอย่าง
\begin{tikzpicture}[font=\small]
  % ---- ซ้าย: รูปสัญลักษณ์ ----
  \node[font=\bfseries] at (-2.7,2.0) {รูปสัญลักษณ์};
  \node[font=\itshape] (p1) at (-2.7,1.15) {$p \rightarrow q$};
  \node[font=\itshape] (p2) at (-2.7,0.45) {$p$};
  \draw[line width=1pt] (-3.9,0.05) -- (-1.5,0.05);
  \node[font=\itshape] at (-2.7,-0.6) {$\therefore\ q$};
  % ---- ขวา: ตัวอย่างรูปธรรม (รูปแบบเดียวกัน) ----
  \node[font=\bfseries] at (2.9,2.0) {ตัวอย่าง};
  \node (e1) at (2.9,1.15) {ถ้าฝนตก แล้วถนนเปียก};
  \node (e2) at (2.9,0.45) {ฝนตก};
  \draw[line width=1pt] (0.5,0.05) -- (5.3,0.05);
  \node at (2.9,-0.6) {$\therefore\ $ ถนนเปียก};
  % เส้นแบ่ง
  \draw[gray,dashed] (-0.1,2.1) -- (-0.1,-0.7);
  % ข้อความสรุปกฎเป็นภาษาง่าย
  \node[draw=gray,fill=cellgray,rounded corners=3pt,align=center,inner sep=2.6mm,
        text width=11.5cm,anchor=north] at (0.4,-1.25)
        {\textbf{กฎ Modus Ponens:} เมื่อ ``$p \rightarrow q$'' และ ``$p$'' เป็นจริงทั้งคู่
         จึงสรุปได้ว่า ``$q$'' เป็นจริง};
\end{tikzpicture}
